
\documentclass{article}
\usepackage{colortbl}
\usepackage{makecell}
\usepackage{multirow}
\usepackage{supertabular}

\begin{document}

\newcounter{utterance}

\centering \large Interaction Transcript for game `cladder', experiment `full\_v1.5\_default', episode 2677 with refineonly{-}llama{-}3.1{-}8b{-}instruct.
\vspace{24pt}

{ \footnotesize  \setcounter{utterance}{1}
\setlength{\tabcolsep}{0pt}
\begin{supertabular}{c@{$\;$}|p{.15\linewidth}@{}p{.15\linewidth}p{.15\linewidth}p{.15\linewidth}p{.15\linewidth}p{.15\linewidth}}
    \# & \multicolumn{2}{c}{Player} && \multicolumn{2}{c}{Game Master} \\
    \hline

    \theutterance \stepcounter{utterance}  
    & & & \multicolumn{4}{p{0.6\linewidth}}{
        \cellcolor[rgb]{0.9,0.9,0.9}{
            \makecell[{{p{\linewidth}}}]{
                \texttt{\tiny{[P1$\langle$GM]}}
                \texttt{You are an expert in causal inference. The following question is not a typical commonsense query, but rather a meticulously designed question created by a professor specializing in causal inference, intended to assess the students' mastery of the course content.} \\
\texttt{Imagine a self{-}contained, hypothetical world with only the following conditions, and without any unmentioned factors or causal relationships: Gene has a direct effect on smoking and lung cancer. Pollution has a direct effect on smoking and lung cancer. Smoking has a direct effect on lung cancer. Pollution is unobserved. Method 1: We look at how gene correlates with lung cancer case by case according to smoking. Method 2: We look directly at how gene correlates with lung cancer in general. To understand how gene affects lung cancer, is it more correct to use the Method 1 than Method 2?} \\
\texttt{Start your answer with 'Yes' or 'No', followed by additional reasoning or evidence to support your explanation.} \\
            }
        }
    }
     \\ \\

    \theutterance \stepcounter{utterance}  
    & & \multicolumn{4}{p{0.6\linewidth}}{
        \cellcolor[rgb]{0.95,0.95,0.95}{
            \makecell[{{p{\linewidth}}}]{
                \texttt{\tiny{[P1$\leftrightarrow$P]}}
                \texttt{{[}DRAFT{]}} \\
\texttt{No, it is not more correct to use Method 1 than Method 2.} \\
\\ 
\texttt{In this hypothetical world, we have a complex causal structure involving Gene, Smoking, and Lung Cancer, with Pollution being an unobserved confounder. } \\
\\ 
\texttt{Method 1, which looks at how Gene correlates with Lung Cancer case by case according to Smoking, is essentially a form of stratification or matching on Smoking. However, this approach does not account for the unobserved confounder Pollution, which has a direct effect on both Smoking and Lung Cancer. } \\
\\ 
\texttt{In the presence of an unobserved confounder, stratification or matching on a related variable (in this case, Smoking) can lead to biased estimates of the causal effect of Gene on Lung Cancer. This is because the strata or matched groups may still differ in terms of the unobserved confounder, leading to confounding bias.} \\
\\ 
\texttt{On the other hand, Method 2, which looks directly at how Gene correlates with Lung Cancer in general, is a more straightforward approach that does not involve stratification or matching. However, this approach is also subject to confounding bias due to the unobserved confounder Pollution.} \\
\\ 
\texttt{To address this issue, we can use a technique called instrumental variable analysis (IVA), which can help to identify the causal effect of Gene on Lung Cancer by leveraging the relationship between Gene and Pollution. However, in the absence of additional information or assumptions, neither Method 1 nor Method 2 provides a reliable estimate of the causal effect of Gene on Lung Cancer.} \\
\\ 
\texttt{In general, when dealing with complex causal structures involving unobserved confounders, it is essential to use techniques that can account for these confounders, such as IVA, regression discontinuity design, or other forms of causal inference that can help to identify the causal effect of interest.} \\
            }
        }
    }
    & \\ \\

    \theutterance \stepcounter{utterance}  
    & \multicolumn{4}{p{0.6\linewidth}}{
        \cellcolor[rgb]{0.561,0.737,0.561}{
            \makecell[{{p{\linewidth}}}]{
                \texttt{\tiny{[P1$\rangle$GM]}}
                \texttt{No, it is not more correct to use Method 1 than Method 2.} \\
\\ 
\texttt{In this hypothetical world, we have a complex causal structure involving Gene, Smoking, and Lung Cancer, with Pollution being an unobserved confounder. } \\
\\ 
\texttt{Method 1, which looks at how Gene correlates with Lung Cancer case by case according to Smoking, is essentially a form of stratification or matching on Smoking. However, this approach does not account for the unobserved confounder Pollution, which has a direct effect on both Smoking and Lung Cancer. } \\
\\ 
\texttt{In the presence of an unobserved confounder, stratification or matching on a related variable (in this case, Smoking) can lead to biased estimates of the causal effect of Gene on Lung Cancer. This is because the strata or matched groups may still differ in terms of the unobserved confounder, leading to confounding bias.} \\
\\ 
\texttt{On the other hand, Method 2, which looks directly at how Gene correlates with Lung Cancer in general, is a more straightforward approach that does not involve stratification or matching. However, this approach is also subject to confounding bias due to the unobserved confounder Pollution.} \\
\\ 
\texttt{To address this issue, we can use a technique called instrumental variable analysis (IVA), which can help to identify the causal effect of Gene on Lung Cancer by leveraging the relationship between Gene and Pollution. However, in the absence of additional information or assumptions, neither Method 1 nor Method 2 provides a reliable estimate of the causal effect of Gene on Lung Cancer.} \\
\\ 
\texttt{In general, when dealing with complex causal structures involving unobserved confounders, it is essential to use techniques that can account for these confounders, such as IVA, regression discontinuity design, or other forms of causal inference that can help to identify the causal effect of interest.} \\
            }
        }
    }
    & & \\ \\

    \theutterance \stepcounter{utterance}  
    & & & \multicolumn{2}{p{0.3\linewidth}}{
        \cellcolor[rgb]{0.9,0.9,0.9}{
            \makecell[{{p{\linewidth}}}]{
                \texttt{\tiny{[GM$|$GM]}}
                \texttt{No} \\
            }
        }
    }
    & & \\ \\

    \theutterance \stepcounter{utterance}  
    & & & \multicolumn{2}{p{0.3\linewidth}}{
        \cellcolor[rgb]{0.9,0.9,0.9}{
            \makecell[{{p{\linewidth}}}]{
                \texttt{\tiny{[GM$|$GM]}}
                \texttt{no} \\
            }
        }
    }
    & & \\ \\

    \theutterance \stepcounter{utterance}  
    & & & \multicolumn{2}{p{0.3\linewidth}}{
        \cellcolor[rgb]{0.9,0.9,0.9}{
            \makecell[{{p{\linewidth}}}]{
                \texttt{\tiny{[GM$|$GM]}}
                \texttt{game\_result = WIN} \\
            }
        }
    }
    & & \\ \\

\end{supertabular}
}

\end{document}
